\section{Conclusion and Future Work}\label{sec:conclusion}

This work introduces a leakage-safe, mechanism-aware protocol for single-sensor imputation in industrial settings and validates it on sulfate production data with extreme target sparsity (80.29\% missingness). The main contribution is methodological: a two-phase framework that avoids optimistic bias through split-first, train-only preprocessing and mechanism-aware controlled validation aligned with deployment.

\paragraph{Contributions and empirical validation.}
The protocol has two pillars. First, a \emph{leakage-safe} discipline (split-first, train-only preprocessing) and an \emph{adaptive head-skip} rule for sparse temporal training ($s^*=0.361$ here) eliminate manual curation and leakage. Second, \emph{mechanism-aware validation} injects missingness only in the test target under MCAR/MAR/MNAR and scores only on masked indices, matching deployment.

Empirically, with strict temporal blocked splitting, CDAE reaches $R^2=\phaseOneRtwo$ in Phase~1 (MAE=\phaseOneMAE, RMSE=\phaseOneRMSE; $n=\phaseOneN$). In Phase~2 (80\% masking), performance remains high under MCAR/MAR ($R^2 \ge 0.93$) and drops under MNAR ($R^2=\mnarRtwo$), as expected under value-dependent masking plus temporal shift. CDAE consistently outperforms KNN-Imputer (about 1.6--2.1x lower MAE, with the largest margin under MNAR, where KNN falls to \knnMnarRtwo). The ablation shows that $\lambda=0$ outperforms $\lambda=1$ on global metrics under temporal evaluation (up to $\ablationLambdaZeroMcarRtwo$), indicating that the deployment-aligned conditional formulation---not loss weighting---is the primary driver of overall accuracy. However, band-restricted analysis (\S\ref{sec:bias-analysis}) reveals a regime-specific trade-off: $\lambda{=}1$ reduces low-flow bias by ${\approx}\,0.4$~EU compared to $\lambda{=}0$ (paired CI excludes zero), acting as implicit importance weighting for the sparse regime at the cost of global MAE.

\paragraph{Realistic assessment.}
Mechanism-aware analysis reveals a density-driven conservative overestimation in the low-flow regime ($[0,20]$~EU; \S\ref{sec:bias-analysis}). The downstream alert pipeline still shows operational value (PR-AUC=\downstreamPRAUC\ vs \downstreamPRAUCBaseline), reducing false positives while preserving sensitivity. Compared with random-split studies, these temporally blocked results are more conservative and therefore more realistic: they expose both strengths (consistent baseline gains, secondary-target generalization with $R^2 \ge 0.979$) and limits (MNAR sensitivity and low-flow bias) under the intended single-sensor outage scenario.

\paragraph{Future directions.}
Future work should further study the interaction between loss weighting and temporal shift, including adaptive weighting schemes for nonstationary operation. Broader validation across processes, horizons, and failure modes is also needed to quantify generalizability beyond sulfate production. Finally, low-flow bias mitigation (oversampling, band-weighted losses, quantile objectives, or regime-specific heads) is a direct path to improve reliability during start-up and shutdown.
